% ============================================================================
% LEHRERHINWEISE: Wiederholung + Stundenleistung Mi casa
% ============================================================================
% Sequenz 3, Block d: Kurzplan, Differenzierung, Losungen
% ============================================================================

\documentclass[11pt, a4paper]{article}

% ============================================================================
% SW-MODUS TOGGLE
% ============================================================================
\newif\ifswmodus
\swmodusfalse

% ============================================================================
% PAKETE
% ============================================================================

\usepackage[utf8]{inputenc}
\usepackage[T1]{fontenc}
\usepackage[ngerman]{babel}
\usepackage{helvet}
\renewcommand{\familydefault}{\sfdefault}

\usepackage[margin=2cm]{geometry}
\usepackage{enumitem}
\usepackage{tabularx}
\usepackage{booktabs}
\usepackage{longtable}

\usepackage{xcolor}
\usepackage{amssymb}
\usepackage{tcolorbox}
\tcbuselibrary{skins}

\usepackage{fancyhdr}
\usepackage{lastpage}
\usepackage{needspace}

% ============================================================================
% FARBEN
% ============================================================================

\definecolor{spanisch-color}{HTML}{c41e3a}
\definecolor{grau-color}{HTML}{888888}
\definecolor{hellgrau-color}{HTML}{f5f5f5}
\definecolor{basis-color}{HTML}{2a7a4b}
\definecolor{standard-color}{HTML}{2c5aa0}
\definecolor{erweiterung-color}{HTML}{7b1fa2}
\definecolor{loesung-color}{HTML}{c62828}

\definecolor{sw-dark}{HTML}{000000}
\definecolor{sw-gray}{HTML}{666666}
\definecolor{sw-white}{HTML}{ffffff}

\ifswmodus
    \colorlet{spanisch}{sw-dark}
    \colorlet{grau}{sw-gray}
    \colorlet{hellgrau}{sw-white}
    \colorlet{basis}{sw-dark}
    \colorlet{standard}{sw-dark}
    \colorlet{erweiterung}{sw-dark}
    \colorlet{loesung}{sw-dark}
\else
    \colorlet{spanisch}{spanisch-color}
    \colorlet{grau}{grau-color}
    \colorlet{hellgrau}{hellgrau-color}
    \colorlet{basis}{basis-color}
    \colorlet{standard}{standard-color}
    \colorlet{erweiterung}{erweiterung-color}
    \colorlet{loesung}{loesung-color}
\fi

\newcommand{\fachfarbe}{spanisch}

% ============================================================================
% VARIABLEN
% ============================================================================

\newcommand{\thema}{Wiederholung + Stundenleistung}
\newcommand{\sequenz}{03}
\newcommand{\block}{d}
\newcommand{\fach}{Spanisch}
\newcommand{\klasse}{7}
\newcommand{\dauer}{90 Min. (Doppelstunde)}
\newcommand{\lerngruppengroesse}{24}
\newcommand{\datum}{\today}

% ============================================================================
% KOPF- UND FUSSZEILE
% ============================================================================

\pagestyle{fancy}
\fancyhf{}
\renewcommand{\headrulewidth}{0.4pt}
\renewcommand{\footrulewidth}{0.4pt}

\lhead{\fach{} -- Klasse \klasse}
\chead{\textbf{LH-\sequenz\block{} (Lehrerhinweise)}}
\rhead{\datum}

\lfoot{}
\cfoot{}
\rfoot{Seite \thepage\ von \pageref{LastPage}}

% ============================================================================
% BEFEHLE
% ============================================================================

\newcommand{\B}{\textcolor{basis}{\textbf{[B]}}}
\newcommand{\Std}{\textcolor{standard}{\textbf{[S]}}}
\newcommand{\E}{\textcolor{erweiterung}{\textbf{[E]}}}

\newcommand{\loesung}[1]{\textcolor{loesung}{\textbf{#1}}}
\newcommand{\es}[1]{\textcolor{\fachfarbe}{\textbf{#1}}}

\newtcolorbox{kurzplan}{
    colback=hellgrau,
    colframe=\fachfarbe,
    colbacktitle=white,
    coltitle=\fachfarbe,
    boxrule=1pt,
    arc=0pt,
    fonttitle=\bfseries,
    attach boxed title to top left={yshift=-2mm, xshift=5mm},
    boxed title style={boxrule=0pt, colback=white},
    title=Kurzplan
}

\newtcolorbox{differenzierung}{
    colback=white,
    colframe=grau,
    colbacktitle=white,
    coltitle=grau,
    boxrule=0.5pt,
    arc=0pt,
    fonttitle=\bfseries,
    attach boxed title to top left={yshift=-2mm, xshift=5mm},
    boxed title style={boxrule=0pt, colback=white},
    title=Differenzierung
}

\newtcolorbox{fehlerbox}{
    colback=white,
    colframe=loesung,
    colbacktitle=white,
    coltitle=loesung,
    boxrule=0.5pt,
    arc=0pt,
    fonttitle=\bfseries,
    attach boxed title to top left={yshift=-2mm, xshift=5mm},
    boxed title style={boxrule=0pt, colback=white},
    title=Haufige Fehler
}

\newtcolorbox{materialbox}{
    colback=white,
    colframe=\fachfarbe,
    colbacktitle=white,
    coltitle=\fachfarbe,
    boxrule=0.5pt,
    arc=0pt,
    fonttitle=\bfseries,
    attach boxed title to top left={yshift=-2mm, xshift=5mm},
    boxed title style={boxrule=0pt, colback=white},
    title=Material-Checkliste
}

% ============================================================================
% DOKUMENT
% ============================================================================

\begin{document}

% ============================================================================
% KOPFBEREICH
% ============================================================================

\begin{center}
    {\Large\bfseries\textcolor{\fachfarbe}{Lehrerhinweise: \thema}}\\[0.3em]
    {\normalsize LH-\sequenz\block{} | \dauer}
\end{center}

\vspace{10pt}

% ============================================================================
% 1. KURZPLAN
% ============================================================================

\section*{1. Stundenverlauf (Kurzplan)}

\textbf{1. Stunde: Wiederholung (45 Min.)}

\begin{tabularx}{\textwidth}{|c|l|X|l|}
\hline
\textbf{Zeit} & \textbf{Phase} & \textbf{Inhalt} & \textbf{Material} \\
\hline
0--5 & Einstieg & Begrußung, Ankun\-digung Ablauf & Tafel \\
\hline
5--10 & Aktivierung & Schnelles Vokabelquiz: Raume, Mobel & Bilder \\
\hline
10--25 & Wiederholung & AB-03d oder LP-03d bearbeiten & AB/LP \\
\hline
25--35 & Partnerubung & Zimmer beschreiben (Aufgabe 4) & -- \\
\hline
35--40 & Sicherung & Fragen klaren, Tipps fur SL & Tafel \\
\hline
40--45 & Uberleitung & Pause, Vorbereitung SL & -- \\
\hline
\end{tabularx}

\vspace{1em}

\textbf{2. Stunde: Stundenleistung (45 Min.)}

\begin{tabularx}{\textwidth}{|c|l|X|l|}
\hline
\textbf{Zeit} & \textbf{Phase} & \textbf{Inhalt} & \textbf{Material} \\
\hline
0--5 & Vorbereitung & SL austeilen, Regeln, Zeit: 25 Min. & SL-03d \\
\hline
5--30 & Durchfuhrung & SuS bearbeiten SL (Einzelarbeit) & SL-03d \\
\hline
30--35 & Einsammeln & SL einsammeln, Ruhe bewahren & -- \\
\hline
35--40 & Reflexion & ,,Was war leicht/schwer?'' & -- \\
\hline
40--45 & Abschluss & Ausblick nachste Sequenz, HA & Tafel \\
\hline
\end{tabularx}

% ============================================================================
% 2. DIFFERENZIERUNG
% ============================================================================

\section*{2. Differenzierung}

\begin{differenzierung}
\textbf{Legende:} \B{} Basis (BR) | \Std{} Standard (MR) | \E{} Erweiterung (GY)

\vspace{0.5em}

\textbf{WICHTIG:} Diese Marker sind NUR fur die Lehrkraft. Auf Schulermaterial erscheinen KEINE Niveaumarkierungen!
\end{differenzierung}

\vspace{1em}

\begin{tabularx}{\textwidth}{|l|X|X|X|}
\hline
& \B{} \textbf{Basis} & \Std{} \textbf{Standard} & \E{} \textbf{Erweiterung} \\
\hline
\textbf{Wiederholung} & Mit Vokabelliste & Selbststandig & + Zusatzfragen \\
\hline
\textbf{Aufg. 1-3 (SL)} & Mit Wortliste & Ohne Hilfe & Schneller fertig \\
\hline
\textbf{Aufg. 4 (SL)} & 3 Satze genugen & 4 Satze & 6+ Satze, komplex \\
\hline
\textbf{Zeitzugabe} & +5 Min. (LRS) & -- & -- \\
\hline
\end{tabularx}

\vspace{1em}

\textbf{Konkrete Hinweise:}
\begin{itemize}
    \item \B{} SuS mit LRS: Zeitzugabe +5 Min., großere Schrift (SL separat drucken)
    \item \B{} DaZ-SuS: Vor der SL Schlusselworter klaren
    \item \E{} Leistungsstarke: Aufgabe 4 mit Zusatzanforderung (,,Beschreibe 2 Zimmer'')
\end{itemize}

% ============================================================================
% 3. HAUFIGE FEHLER
% ============================================================================

\section*{3. Haufige Fehler}

\begin{fehlerbox}
\begin{tabularx}{\textwidth}{|l|X|X|}
\hline
\textbf{Fehler} & \textbf{Beispiel} & \textbf{Reaktion} \\
\hline
Genus bei Mobeln & *el cama statt \es{la cama} & Merkregel: -a $\rightarrow$ feminin \\
\hline
Adjektivanpassung & *mesa rojo statt \es{mesa roja} & Visualisierung: Endungen markieren \\
\hline
hay mit Artikel & *hay la mesa statt \es{hay una mesa} & Kontrastieren: hay una vs. la mesa es \\
\hline
Plural bei Farben & *verdes statt \es{verdes} & Plural -es bei Konsonanten \\
\hline
\end{tabularx}
\end{fehlerbox}

% ============================================================================
% 4. MATERIAL-CHECKLISTE
% ============================================================================

\section*{4. Material-Checkliste}

\begin{materialbox}
\begin{itemize}[label=$\square$]
    \item AB-03d.pdf (\lerngruppengroesse{} Kopien)
    \item ML-03d.pdf (1x fur Lehrkraft)
    \item SL-03d.pdf (\lerngruppengroesse{} Kopien + Reserve)
    \item SL-03d-ML.pdf (1x fur Lehrkraft)
    \item Lehrwerk: Qué pasa 1, S. 36-45
    \item Tafelmarker / Kreide
    \item Optional: LP-03d.html (Tablets/PCs)
    \item Optional: Bildkarten Raume/Mobel
    \item Uhr / Timer fur 25 Min.
\end{itemize}
\end{materialbox}

% ============================================================================
% 5. LOSUNGEN AB-03d
% ============================================================================

\section*{5. Losungen AB-03d (Wiederholung)}

\textbf{Merkkasten:}
\begin{tabular}{ll}
\es{la cocina} & $\rightarrow$ \loesung{die Kuche} \\
\es{el salon} & $\rightarrow$ \loesung{das Wohnzimmer} \\
\es{el dormitorio} & $\rightarrow$ \loesung{das Schlafzimmer} \\
\es{el bano} & $\rightarrow$ \loesung{das Badezimmer} \\
\es{el jardin} & $\rightarrow$ \loesung{der Garten} \\
\end{tabular}

\vspace{1em}

\textbf{Aufgabe 1: Zuordnung} (4 Punkte)
\begin{tabular}{ll}
la cocina & $\rightarrow$ \loesung{die Kuche} \\
el salon & $\rightarrow$ \loesung{das Wohnzimmer} \\
el dormitorio & $\rightarrow$ \loesung{das Schlafzimmer} \\
el jardin & $\rightarrow$ \loesung{der Garten} \\
\end{tabular}

\vspace{1em}

\textbf{Aufgabe 2: hay + Mobel} (4 Punkte)
\begin{tabular}{ll}
a) & \loesung{hay una cama} \\
b) & \loesung{hay un sofa} \\
c) & \loesung{hay una mesa} \\
d) & \loesung{hay un armario} \\
\end{tabular}

\vspace{1em}

\textbf{Aufgabe 3: Farbadjektive} (6 Punkte)
\begin{tabular}{ll}
a) & \loesung{roja} \\
b) & \loesung{azul} \\
c) & \loesung{amarilla} \\
d) & \loesung{blanco} \\
e) & \loesung{verdes} \\
f) & \loesung{negros} \\
\end{tabular}

\vspace{1em}

\textbf{Aufgabe 4: Freie Produktion} (6 Punkte)

Bewertungskriterien:
\begin{itemize}
    \item Mindestens 4 Satze: 2P
    \item Korrekter Gebrauch \textit{hay}: 2P
    \item Adjektive angepasst: 1P
    \item Sprachliche Korrektheit: 1P
\end{itemize}

% ============================================================================
% 6. LOSUNGEN SL-03d (STUNDENLEISTUNG)
% ============================================================================

\section*{6. Losungen SL-03d (Stundenleistung)}

\textbf{Aufgabe 1: Raume zuordnen} (4 Punkte)
\begin{tabular}{ll}
1. la cocina & $\rightarrow$ \loesung{c) die Kuche} \\
2. el dormitorio & $\rightarrow$ \loesung{a) das Schlafzimmer} \\
3. el salon & $\rightarrow$ \loesung{d) das Wohnzimmer} \\
4. el bano & $\rightarrow$ \loesung{b) das Badezimmer} \\
\end{tabular}

\vspace{1em}

\textbf{Aufgabe 2: hay + Mobel} (4 Punkte)
\begin{tabular}{ll}
a) & \loesung{hay una mesa} \\
b) & \loesung{hay una cama} \\
c) & \loesung{hay un sofa} \\
d) & \loesung{hay un armario} \\
\end{tabular}

\vspace{1em}

\textbf{Aufgabe 3: Farben + Adjektivanpassung} (6 Punkte)
\begin{tabular}{ll}
a) el sofa es & \loesung{verde} \\
b) la silla es & \loesung{roja} \\
c) las mesas son & \loesung{blancas} \\
d) los armarios son & \loesung{marrones} \\
e) la cama es & \loesung{azul} \\
f) el jardin es & \loesung{grande} \\
\end{tabular}

\vspace{1em}

\textbf{Aufgabe 4: Zimmer beschreiben} (6 Punkte)

Bewertungskriterien:
\begin{itemize}
    \item Mindestens 4 Satze: 2P
    \item Korrekter Gebrauch \textit{hay}: 2P
    \item Farbadjektive korrekt: 1P
    \item Sprachliche Korrektheit: 1P
\end{itemize}

% ============================================================================
% 7. BEWERTUNGSSKALA
% ============================================================================

\section*{7. Bewertungsskala (14-NP, Sek I)}

\begin{center}
\begin{tabular}{|c|c|c|c|c|c|c|c|c|c|c|c|c|c|c|}
\hline
\textbf{NP} & 14 & 13 & 12 & 11 & 10 & 9 & 8 & 7 & 6 & 5 & 4 & 3 & 2 & 1 \\
\hline
\textbf{\%} & 100 & 96 & 91 & 86 & 80 & 73 & 67 & 60 & 53 & 47 & 40 & 33 & 27 & 20 \\
\hline
\textbf{Punkte} & 20 & 19 & 18 & 17 & 16 & 15 & 13 & 12 & 11 & 9 & 8 & 7 & 5 & 4 \\
\hline
\end{tabular}
\end{center}

% ============================================================================
% 8. HAUSAUFGABE
% ============================================================================

\section*{8. Hausaufgabe}

\begin{itemize}
    \item Vokabeln Sequenz 3 wiederholen (Vorbereitung Klassenarbeit)
    \item Optional: Lernpfad LP-03d.html (digital)
\end{itemize}

% ============================================================================
% 9. REFLEXIONSNOTIZEN
% ============================================================================

\section*{9. Reflexionsnotizen (nach Durchfuhrung)}

\begin{tabularx}{\textwidth}{|l|X|}
\hline
\textbf{Was lief gut?} & \\[2em]
\hline
\textbf{Was lief nicht?} & \\[2em]
\hline
\textbf{Zeitmanagement} & \\[2em]
\hline
\textbf{Fur nachstes Mal} & \\[2em]
\hline
\end{tabularx}

\end{document}
